\section{Conclusioni}

Con questa esperienza abbiamo potuto studiare lo spettro di emissione gamma di alcuni isotopi radioattivi e la loro attività.  
\help{INSERIRE PARTI CONCLUSIVE}
L'analisi delle sorgenti $FS65$ ed $FS66$ di \ce{^60Cs} ha portato a risultati compatibili con quelli attesi sia utilizzando il metodo relativo che quello assoluto. 
\begin{equation}
	A_{FS66\text{rel}} = (36.8 \pm 1.9) \unit{kBq}
\end{equation}	
\begin{equation}
	A_{FS65\text{abs}} = (196 \pm 17) \unit{kBq} \quad A_{FS66\text{abs}} = (38 \pm 3) \unit {kBq}
\end{equation}
Successivamente si è ricavato il coefficiente di assorbimento di massa del piombo che è risultato essere statisticamente compatibile con il valore teorico.
\begin{equation}
	\mu /\rho = (0.1040 \pm 0.0010) \unit{cm^2/g} 
\end{equation}
Infine, ci siamo preoccupati di studiare il picco somma del \ce{^60Cs} in modo da poter determinare se gli eventi nel picco fossero correlati o se avvenissero in maniera indipendente. L'attività trovata in laboratorio risulta essere incompatibile con entrambe le ipotesi. 
\begin{equation}
	A_{FS64\text{today}} = (32.5 \pm 1.7) \unit{kBq} 
\end{equation}
\help{Come risultato di questa misura io ho messo l'attività ma non penso sia corretto, ha più senso mettere i valori di $\Sigma_1$ e $\Sigma_2$?}
Tuttavia il risultato suggerisce che il decadimento sia a cascata visto che i due valori hanno lo stesso ordine di grandezza. 
